% =============================================================================
%  Manual del estudiante · Proyecto final, semana 1 («El hilo»)
%  Razonamiento Automático (prof. Fabián Serrano) · UNAL Manizales
%  Compilar: pdflatex Manual_Semana1.tex (dos veces)
% =============================================================================
\documentclass[11pt,a4paper]{article}
\usepackage[utf8]{inputenc}
\usepackage[T1]{fontenc}
\IfFileExists{spanish.ldf}
  {\usepackage[spanish,es-noshorthands,es-tabla]{babel}}
  {\usepackage[provide=*]{babel}\babelprovide[import=es, main]{spanish}}
\usepackage{mathpazo}
\linespread{1.06}
\usepackage[margin=2.3cm,top=2.4cm,bottom=2.6cm]{geometry}
\usepackage{amsmath,amssymb}
\usepackage{xcolor}
\usepackage{booktabs,array,tabularx}
\usepackage{enumitem}
\usepackage{listings}
\usepackage{titlesec}
\usepackage{fancyhdr}
\usepackage[breakable,skins]{tcolorbox}
\usepackage[hidelinks]{hyperref}

\definecolor{psKey}{RGB}{0,80,160}
\definecolor{psStr}{RGB}{160,60,0}
\definecolor{resalte}{RGB}{180,0,60}
\definecolor{confiable}{RGB}{0,120,60}
\definecolor{noconfiable}{RGB}{200,110,0}
\definecolor{humano}{RGB}{0,80,160}

\setlist{itemsep=2pt,topsep=4pt}
\setlength{\parindent}{0pt}
\setlength{\parskip}{5pt}
\setcounter{secnumdepth}{1}
\titleformat{\section}{\Large\sffamily\bfseries\color{humano}}{\thesection.}{0.5em}{}
\titlespacing*{\section}{0pt}{16pt}{6pt}
\titleformat{\subsection}{\large\sffamily\bfseries}{\thesubsection}{0.5em}{}
\titlespacing*{\subsection}{0pt}{10pt}{3pt}

\pagestyle{fancy}
\fancyhf{}
\fancyhead[L]{\small\sffamily\textcolor{gray}{Manual del estudiante · semana 1}}
\fancyhead[R]{\small\sffamily\textcolor{gray}{\leftmark}}
\fancyfoot[C]{\small\sffamily\thepage}
\renewcommand{\headrulewidth}{0pt}
\renewcommand{\sectionmark}[1]{\markboth{#1}{}}

\tcbset{enhanced, breakable, boxrule=0.6pt, arc=2pt, left=7pt, right=7pt, top=5pt, bottom=5pt,
  fonttitle=\sffamily\bfseries\small, before skip=8pt, after skip=8pt}
\newtcolorbox{idea}{unbreakable, colback=confiable!6, colframe=confiable!75!black, title={La idea en una frase}}
\newtcolorbox{haga}{colback=humano!5, colframe=humano!75!black, title={Hágalo}}
\newtcolorbox{ojo}{unbreakable, colback=resalte!5, colframe=resalte!80!black, title={Ojo}}
\newtcolorbox{falla}[1]{unbreakable, colback=noconfiable!7, colframe=noconfiable!85!black, title={Si sale: #1}}

\newcommand{\cd}[1]{\texttt{\small #1}}
\lstset{basicstyle=\ttfamily\small, columns=fullflexible, keepspaces=true,
  showstringspaces=false, breaklines=true,
  keywordstyle=\color{psKey}\bfseries, commentstyle=\color{gray}\itshape,
  stringstyle=\color{psStr},
  backgroundcolor=\color{black!4}, frame=leftline, rulecolor=\color{psKey!40},
  framesep=6pt, xleftmargin=8pt, aboveskip=0.5em, belowskip=0.5em}

\begin{document}

\begin{center}
  {\LARGE\sffamily\bfseries\color{humano} Manual del estudiante --- Semana 1}\\[0.4em]
  {\large El hilo: el modelo, la configuración y la línea base}\\[0.6em]
  {\sffamily\small\color{gray} Proyecto final de Razonamiento Automático · Prof.\ Fabián Serrano\\
  Asesoría: Carlos Orrego · UNAL Manizales · 5 de noviembre de 2026}
\end{center}

\begin{idea}
  Al terminar esta semana cada uno tiene un modelo de lenguaje gratuito
  respondiendo desde Python y sabe medir qué tan mal demuestra cuando
  nadie lo ayuda. Ese número se llama \textbf{línea base}, y todo el
  proyecto se mide contra él.
\end{idea}

\textbf{Cómo usar este manual.} Síganlo en orden, con la terminal
abierta. Cada caja azul es algo que hay que hacer; cada caja naranja,
un error típico y cómo salir de él. Todo el código está en el repositorio:
\url{https://github.com/cmorregof/repo_proyecto_final}.

% ====================================================================
\section{La idea del proyecto}

Un modelo de lenguaje predice la siguiente palabra más probable. Por eso
escribe con fluidez, pero lo que escribe es \emph{plausible}, no
necesariamente \emph{verdadero}. Cuando afirma algo falso con total
seguridad, se dice que \textbf{alucina}. Isabelle hace lo contrario: no
se deja convencer por el estilo, acepta o rechaza, y dice dónde falla.

El proyecto junta a los dos: \textbf{el modelo propone; Isabelle
decide.} El asistente que van a construir busca en los documentos del
curso (eso es un \emph{RAG}: un examen a libro abierto), le pide al
modelo una demostración, y se la hace verificar a Isabelle antes de
mostrarla. Si Isabelle la rechaza, el error vuelve al modelo para que
corrija.

\begin{center}
\small
\begin{tabular}{@{}lll@{}}
  \toprule
  Jueves & Nombre & Qué se construye \\
  \midrule
  \textbf{1} & \textbf{El hilo} & \textbf{el modelo y la línea base (configuración 1)} \\
  2 & La biblioteca & buscar en los documentos y citar (configuración 2) \\
  3 & La flecha roja & verificar con Isabelle y corregir (configuración 3) \\
  4 & El espejo & evaluar, analizar errores, informe \\
  \bottomrule
\end{tabular}
\end{center}

Se califica \textbf{la mejora} de la configuración 1 a la 3 con el
mismo modelo, no el porcentaje absoluto: un modelo gratuito pequeño
compite en igualdad con uno de pago. Y por encima de todo, la
honestidad: el asistente nunca dice «verificado» sin que Isabelle lo
haya aceptado.

% ====================================================================
\section{Preparar el computador}

\subsection{Lo que necesitan}
\begin{enumerate}
  \item \textbf{Isabelle} (el del curso).
  \item \textbf{Python 3.9 o posterior.} No hay que instalar paquetes.
  \item \textbf{git} (o descargar el repositorio como ZIP).
  \item \textbf{Un modelo de lenguaje gratuito}: Ollama (recomendado) u
    OpenRouter.
\end{enumerate}

\subsection{Descargar el repositorio}
\begin{haga}
\begin{lstlisting}[language=bash]
git clone https://github.com/cmorregof/repo_proyecto_final.git
cd repo_proyecto_final
\end{lstlisting}
Sin git: botón verde \textbf{Code $\to$ Download ZIP} en la página del
repositorio. Si ya lo habían clonado, actualicen con \cd{git pull}.
\end{haga}

\subsection{Decirle a Python dónde está Isabelle}
\begin{haga}
Si escribir \cd{isabelle version} en la terminal no funciona, indiquen la
ruta (ajústenla a su instalación y versión):
\begin{lstlisting}[language=bash]
# Mac
export ISABELLE=/Applications/Isabelle2025-2.app/bin/isabelle
# Linux
export ISABELLE=$HOME/Isabelle2025-2/bin/isabelle
\end{lstlisting}
Hay que repetirlo en cada terminal nueva (o agregarlo a \cd{\textasciitilde/.zshrc}
o \cd{\textasciitilde/.bashrc}).
\end{haga}

\subsection{Camino A: Ollama con Gemma 3 (recomendado)}
Ollama es un programa gratuito que descarga modelos y los corre en su
propio computador: sin límite de uso, sin llave, sin costo y sin
internet. Gemma 3 4B funciona en un portátil de 8\,GB de RAM, sin
tarjeta gráfica (lento, pero funciona).
\begin{haga}
\begin{enumerate}
  \item Instalar Ollama desde \url{https://ollama.com} (Mac:
    también \cd{brew install ollama}).
  \item Descargar el modelo, una sola vez (unos 3\,GB):
\begin{lstlisting}[language=bash]
ollama pull gemma3:4b
\end{lstlisting}
  \item Dejar el servidor corriendo \textbf{en una terminal aparte} que no
    se cierra mientras trabajan:
\begin{lstlisting}[language=bash]
ollama serve
\end{lstlisting}
\end{enumerate}
\end{haga}
\begin{falla}{\texttt{address already in use}}
  Ollama ya está corriendo (por ejemplo, la aplicación). No hay que hacer
  nada.
\end{falla}

\subsection{Camino B: OpenRouter}
Para quien no tenga un computador que aguante el modelo local. Creen una
cuenta en \url{https://openrouter.ai}, generen una llave y expórtenla:
\begin{lstlisting}[language=bash]
export OPENROUTER_API_KEY=sk-or-...
\end{lstlisting}
Luego, en \cd{semana1/taller/modelo.py}, cambien \cd{"proveedor"} a
\cd{"openrouter"} y \cd{"modelo"} a un modelo cuyo nombre termine en
\cd{:free}.
\begin{ojo}
  La llave es una contraseña: \textbf{nunca} la escriban dentro del
  código ni la suban a GitHub. Las cuentas gratuitas tienen un límite
  diario (del orden de 50 solicitudes); una corrida completa del
  proyecto puede pedir 400. Por eso el camino A es el principal.
\end{ojo}

\subsection{Otros modelos}
Cualquier modelo pequeño y abierto sirve: \cd{qwen2.5-coder:7b} (sabe de
código, pide unos 16\,GB), \cd{qwen3:4b}, \cd{phi4-mini},
\cd{llama3.2:3b} (para computadores modestos). Se cambia el nombre en
\cd{CONFIG}. \textbf{Regla:} las tres configuraciones del informe, con el
mismo modelo.

% ====================================================================
\section{El kit del profesor: la carpeta \texttt{Recursos/}}

\begin{center}
\small
\begin{tabular}{@{}ll@{}}
  \toprule
  \cd{verificador.py} & le pregunta a Isabelle si un texto es correcto \\
  \cd{benchmark\_R1.json} & las 58 preguntas del benchmark público \\
  \cd{evaluar.py} & califica sus respuestas (el mismo que usará el profesor) \\
  \cd{ejemplo\_respuestas.json} & el formato de respuestas \\
  \cd{README.md} & instrucciones detalladas: léanlo \\
  \bottomrule
\end{tabular}
\end{center}

\subsection{Probar el verificador}
\begin{haga}
Desde la carpeta \cd{Recursos/}:
\begin{lstlisting}[language=Python]
from verificador import verificar

r = verificar('lemma "p \\<and> q \\<Longrightarrow> q \\<and> p" by auto')
print(r.ok)          # True si Isabelle la acepta
print(r.segundos)    # unos 4 o 5
print(r.errores)     # lista vacia si r.ok
\end{lstlisting}
\end{haga}
El verificador acepta una teoría completa o solo un pedazo, convierte los
símbolos Unicode ($\wedge$, $\longrightarrow$) a la forma que exige
Isabelle (\cd{\textbackslash<and>}\ldots), da los errores con la línea del
texto que ustedes enviaron, corta las pruebas que no terminan y no deja
pasar \cd{sorry}.

\subsection{Cómo es una pregunta y qué se responde}
Cada pregunta del benchmark tiene \cd{id}, \cd{estilo} (\cd{"dn"},
deducción natural, o \cd{"libre"}), \cd{enunciado} y \cd{contexto} (todo
lo que va antes del enunciado en la teoría: cabecera y lemas previos).

La respuesta es \textbf{solo lo que va después del enunciado}, desde
\cd{proof}, \cd{by}, \cd{using} o \cd{apply}. El evaluador arma
\cd{contexto + enunciado + respuesta + end}.

\begin{center}
\small
\begin{tabular}{@{}ll@{}}
  \toprule
  \cd{aceptada} & Isabelle la acepta y respeta el estilo \\
  \cd{estilo} & la acepta, pero usa \cd{auto} donde se pedía deducción natural \\
  \cd{rechazada} & Isabelle no la acepta; el informe trae el error \\
  \cd{trampa} & contiene lemas, axiomas, \cd{sorry} u otros comandos: ni se verifica \\
  \cd{sin respuesta} & falta \\
  \bottomrule
\end{tabular}
\end{center}

% ====================================================================
\section{El taller: la configuración 1}

Todo está en \cd{semana1/taller/}. Son cuatro archivos y dos programas:

\begin{center}
\small
\begin{tabular}{@{}ll@{}}
  \toprule
  \cd{modelo.py} & \cd{CONFIG} y \cd{preguntar()}: hablar con el modelo \\
  \cd{utilidades.py} & \cd{extraer\_prueba()} y \cd{registrar()} \\
  \cd{prompt.py} & \cd{pedir\_mensajes()}: armar la pregunta \\
  \cd{config1.py} & corre el benchmark con solo el modelo \\
  \cd{paso1c.py} & primera conversación de prueba \\
  \bottomrule
\end{tabular}
\end{center}

\subsection{Paso 1: hablar con el modelo}
En \cd{modelo.py}, todo lo que cambia entre un modelo y otro vive en un
diccionario:
\begin{lstlisting}[language=Python]
CONFIG = {
    "proveedor":   "ollama",      # o "openrouter"
    "modelo":      "gemma3:4b",   # o el nombre de otro modelo
    "temperatura": 0.2,           # baja = respuestas mas estables
}
\end{lstlisting}
La función \cd{preguntar(mensajes)} recibe una lista de mensajes y
devuelve el texto de la respuesta. Ollama y OpenRouter usan el mismo
formato, así que el resto del código no cambia. Los mensajes tienen
\textbf{rol}: \cd{system} (quién es el modelo y sus reglas), \cd{user} (la
pregunta) y \cd{assistant} (lo que el modelo ya contestó; lo usaremos la
semana 3).

\begin{haga}
\begin{lstlisting}[language=bash]
cd semana1/taller
python3 paso1c.py
\end{lstlisting}
Debe aparecer una explicación de la regla \cd{mp}.
\end{haga}
\begin{falla}{\texttt{Connection refused}}
  Ollama no está corriendo. Abran otra terminal, escriban
  \cd{ollama serve} y déjenla abierta.
\end{falla}
\begin{falla}{\texttt{No such file or directory}}
  No están en la carpeta correcta. \cd{cd} hasta
  \cd{repo\_proyecto\_final/semana1/taller} y vuelvan a intentar.
\end{falla}

\subsection{Lean la respuesta con lupa}
En el ensayo, Gemma 3 4B respondió a «¿qué dice la regla \cd{mp}?» con
una explicación correcta de modus ponens\ldots{} y luego un «ejemplo en
Isabelle/Isar» con \cd{module Main}, \cd{assert}, \cd{check} e
\cd{infer [mp]}, y unas reglas \cd{and} y \cd{or}. \textbf{Nada de eso
existe en Isabelle}: una teoría empieza con \cd{theory \ldots{} imports
Main begin}, \cd{mp} se usa con \cd{using 1 2 by (rule mp)}, y las reglas
de la conjunción y la disyunción son \cd{conjI}, \cd{conjE},
\cd{disjI1}, \cd{disjE}. Todo dicho con la misma seguridad.

\begin{ojo}
  Eso es una alucinación: la mitad es verdad y la otra mitad es
  «Isabelle imaginario», sin que el tono cambie. Por eso este proyecto
  existe.
\end{ojo}

\subsection{Paso 2: limpiar la respuesta y dejar registro}
Los modelos rodean la prueba de explicaciones y comillas triples, y
\textbf{a veces repiten el enunciado aunque se les pida que no}.
\cd{extraer\_prueba} saca el bloque de código; si empieza con
\cd{lemma}, corta hasta la primera línea que empiece con \cd{proof},
\cd{by}, \cd{using} o \cd{apply}; y quita un \cd{oops} final.

\cd{registrar} agrega una línea JSON por evento al archivo de registro
(modo \cd{"a"}: se agrega, nunca se borra). \textbf{Sin registros no se
califica} (requisito R5).

\subsection{Paso 3: armar la pregunta}
\cd{pedir\_mensajes(p, extra="")} junta el contexto, un espacio
\cd{extra}, el enunciado y la regla de estilo (si es \cd{dn}: solo
\cd{rule}, \cd{erule}, \cd{assumption}\ldots; nada de \cd{auto},
\cd{simp}, \cd{blast}). Hoy \cd{extra} va vacío; la semana 2 ahí entran
los pasajes que encuentre la búsqueda.

\subsection{Paso 4: correr y calificar}
\begin{haga}
Primero con 3 preguntas, para comprobar que todo funciona:
\begin{lstlisting}[language=bash]
python3 config1.py 3
python3 ../../Recursos/evaluar.py ../../Recursos/benchmark_R1.json \
        respuestas_config1.json --paralelo 4 --informe informe_c1.json
\end{lstlisting}
Cuando funcione, las 58: \cd{python3 config1.py} y de nuevo
\cd{evaluar.py}.
\end{haga}
\begin{ojo}
  Con Gemma 3 4B en un portátil sin tarjeta gráfica, las 58 tardaron unos
  \textbf{22 minutos} en el ensayo: la mayoría tarda segundos, pero
  algunas respuestas se van en bucle y tardan hasta 3 minutos. Déjenlo
  corriendo y hagan otra cosa.
\end{ojo}

\subsection{Lo que pasó en el ensayo: 44 «trampas»}
Con una primera versión de \cd{extraer\_prueba} que no quitaba el
enunciado repetido, las 58 salieron: \textbf{0 aceptadas, 0 de estilo, 14
rechazadas y 44 \cd{trampa}}. Gemma no hacía trampa a propósito:
desobedecía la instrucción «sin repetir el enunciado» y copiaba el
\cd{lemma}, así que el evaluador ni la mandaba a Isabelle.

\textbf{La lección:} una instrucción en el prompt no es una garantía. Lo
que el modelo \emph{debe} cumplir lo verifica el programa. Limpiar el
formato no es hacer trampa (no se cambia la prueba, solo se quita lo que
el evaluador ya pone), pero hay que declararlo en el informe.

Errores reales de Gemma que Isabelle atrapó en el ensayo:
\begin{itemize}
  \item \textbf{Ejercicio 11:} inventa la regla
    \cd{implicacion\_con\_consecuencia}; Isabelle: \cd{Undefined fact}.
  \item \textbf{Ejercicio 12:} inventa la regla \cd{imply}.
  \item \textbf{Ejercicio 1:} \cd{using 1, 2 by rule} falla solo por la
    coma; sin ella, Isabelle la acepta.
  \item \textbf{Ejercicio 1, otra variante:} \cd{using 2 1 by (rule mp)}
    falla con \cd{Failed to apply initial proof method}, porque \cd{mp}
    espera primero la implicación y luego la premisa.
\end{itemize}

\textbf{El número que obtengan con la limpieza corregida es su línea
base. Anótenlo.}

% ====================================================================
\section{Temas de formalización}

Cada grupo (máximo tres) toma un tema distinto. Los ejercicios de cada
grupo serán el examen de los asistentes de los demás: al menos 8, dos de
ellos \cd{[***]}.

\begin{center}
\small
\begin{tabularx}{\linewidth}{@{}>{\raggedright\arraybackslash}p{3.4cm}Xc@{}}
  \toprule
  Rama & Teorema central & Nivel \\
  \midrule
  Codificación & \cd{decodificar (codificar xs) = xs} (longitud de racha) & $\star$ \\
  Criptografía clásica & César o Vigenère: descifrar invierte a cifrar & $\star$ \\
  Algoritmos & Ordenamiento por inserción: ordena y es permutación & $\star\star$ \\
  Estructuras de datos & Árboles de búsqueda: la inserción preserva el orden & $\star\star$ \\
  Compiladores & Expresiones a máquina de pila: \cd{exec (compilar e) [] = [eval e]} & $\star\star$ \\
  Lenguajes formales & Palíndromos o paréntesis balanceados & $\star\star$ \\
  Lógica computacional & La forma normal negativa preserva el significado & $\star\star$ \\
  Arquitectura & Sumador binario sobre listas de bits & $\star\star\star$ \\
  Grafos & Alcanzabilidad como cierre reflexivo-transitivo & $\star\star\star$ \\
  Lenguajes de programación & Tipado de expresiones: preservación & $\star\star\star$ \\
  Teoría de la computación & Autómatas: el producto reconoce la intersección & $\star\star\star$ \\
  Software (entusiastas) & Verificar un programa imperativo (ver \emph{Concrete Semantics}) & $\star\star\star$ \\
  \bottomrule
\end{tabularx}
\end{center}

% ====================================================================
\section{Tarea para el próximo jueves}

\begin{haga}
\begin{enumerate}
  \item Configuración 1 completa sobre las 58: \cd{respuestas\_config1.json},
    \cd{registros/config1.jsonl} e \cd{informe\_c1.json}.
  \item Leer diez registros a mano y clasificar los errores del modelo:
    sintaxis, reglas inventadas, orden de los hechos, \cd{auto}
    prohibido, enunciado repetido\ldots{} Es material del informe.
  \item Tema de formalización: tres opciones en orden de preferencia, y
    los dos primeros enunciados escritos en Isabelle.
\end{enumerate}
\end{haga}

\begin{center}
  \large\color{resalte}\bfseries La verdad no la pone el modelo más caro: la pone Isabelle.
\end{center}

\end{document}
